% Generated draft: requires mathematical and source review.
\documentclass[11pt]{article}
\usepackage[margin=25mm]{geometry}
\usepackage{fontspec}
\setmainfont{Latin Modern Roman}
\usepackage{amsmath,amssymb,mathtools,unicode-math}
\setmathfont{Latin Modern Math}
\usepackage{enumitem,needspace,xurl,hyperref}
\setlength{\parindent}{0pt}
\setlength{\parskip}{5pt}
\setlength{\emergencystretch}{3em}
\newcommand{\N}{\mathbb{N}}
\newcommand{\Z}{\mathbb{Z}}
\newcommand{\Q}{\mathbb{Q}}
\newcommand{\R}{\mathbb{R}}
\newcommand{\C}{\mathbb{C}}
\newcommand{\F}{\mathbb{F}}
\newcommand{\A}{\mathbb{A}}
\newcommand{\PP}{\mathbb{P}}
\newcommand{\E}{\mathbb{E}}
\newcommand{\eps}{\varepsilon}
\newcommand{\dd}{\,\mathrm{d}}
\begin{document}
\section*{UnsolvedMath ID: 30001266}
\textbf{Source problem code:} catalogue entry
\subsection{Definitions and mathematical statement}
Let \(d\) denote the manifold dimension; the supplied record does not specify its permitted range. Let \(\gamma_d\to B\operatorname{SO}(d)\) be the universal oriented real \(d\)-plane bundle, and let
\[
MTSO(d)=\operatorname{Th}(-\gamma_d),\qquad X_d=\Omega^\infty_0MTSO(d),
\]
where \(X_d\) is the path component of the basepoint in the infinite loop space of this Thom spectrum. All cohomology below has coefficients in \(\Q\). Write
\[
L(\gamma_d)=1+L_1(\gamma_d)+L_2(\gamma_d)+\cdots,
\qquad L_k(\gamma_d)\in H^{4k}(B\operatorname{SO}(d);\Q),
\]
for the Hirzebruch \(L\)-class. Under the usual Thom and cohomology-suspension interpretation of the source notation, put
\[
\mathrm{th}:H^{4k}(B\operatorname{SO}(d);\Q)\longrightarrow H^{4k-d}(MTSO(d);\Q),
\qquad
s:H^*(MTSO(d);\Q)\longrightarrow H^*(X_d;\Q),
\]
and write \(s\,\mathrm{th}(L_k)\) for \(s(\mathrm{th}(L_k(\gamma_d)))\). These maps explain the notation conventionally; their conventions are not defined in the supplied record.

A homotopy-commutative topological monoid \(M\) has a multiplication \(\mu:M\times M\to M\) homotopic to its composite with interchange of the factors. Its path components form a monoid \(\pi_0(M)\). A natural interpretation of the superscript in the source is invariance under translation by components: for \(a\in\pi_0(M)\), choose \(m_a\in M\) representing \(a\), let \(T_a(x)=\mu(m_a,x)\), and set
\[
H^*(M;\Q)^{\pi_0(M)}
=\{u\in H^*(M;\Q):T_a^*u=u\text{ for every }a\in\pi_0(M)\}.
\]
This interpretation must be checked against the original report, since the supplied record does not specify the action.

The supplied existence question asks for such a monoid \(M\), constructed from classifying spaces of diffeomorphism groups of oriented \(d\)-manifolds with boundary, with multiplication induced by gluing, and for the intended geometric comparison homomorphism
\[
\Phi_d:H^*(X_d;\Q)\longrightarrow H^*(M;\Q)^{\pi_0(M)}
\]
to be surjective with
\[
\ker\Phi_d=
\begin{cases}
0,&d\text{ even},\\[2pt]
\bigl\langle s\,\mathrm{th}(L_k):4k\ge d\bigr\rangle,&d\text{ odd}.
\end{cases}
\]
Here the angle brackets denote the ideal in \(H^*(X_d;\Q)\) generated by the displayed classes. The inequality \(4k\ge d\), including its equality case, is retained exactly from the supplied formulation. The original statement attributes the conjectured kernel to J. Ebert.

The record does not specify the allowable manifolds, compactness or connectedness assumptions, boundary identifications, boundary conditions on diffeomorphisms, the precise gluing construction, or the construction of \(\Phi_d\). Consequently, this is a conditional formulation of the recorded target, not a fully specified geometric existence statement. The intended comparison map cannot be replaced without justification by an arbitrary algebra homomorphism.
\subsection{Short English statement}
The question asks whether gluing oriented \(d\)-manifolds with boundary can produce a homotopy-commutative monoid whose translation-invariant rational cohomology is a quotient of \(H^*(\Omega^\infty_0MTSO(d);\Q)\), with the stated parity-dependent kernel attributed to J. Ebert. The manifold class, gluing data, comparison map, and dimension range remain unspecified. The supplied dated assessment labels the problem partially solved and says that Madsen--Weiss settles the surface case; this assessment has not been independently verified.
\subsection{Sources}
\begin{itemize}
\item Supplied local UnsolvedMath registry record, permanent ID \(30001266\), source problem code \texttt{catalogue entry}, title \emph{Stable Cohomology of Higher-Dimensional Diffeomorphism Groups}; collection: Oberwolfach Reports Open Problems. The record gives proposed year 2009 and no named proposer. Its original statement attributes the kernel conjecture to J. Ebert.
\item Supplied source citation: \emph{Manifold Perspectives} (2009), Oberwolfach Reports, pp.\ 1487--1546, DOI \texttt{10.4171/owr/2009/27}. Source URL and literature-source link labelled ``original OWR report'': \url{https://doi.org/10.4171/owr/2009/27}. The supplied background describes Oberwolfach Reports as documenting mathematical workshops and their open-problem discussions.
\item Supplied literature-source link labelled ``Madsen--Weiss theorem'': \url{https://annals.math.princeton.edu/2007/165-3/p04}.
\item Supplied literature-source link labelled ``higher-dimensional MTW results'': \url{https://arxiv.org/abs/0902.4719}.
\item Literature assessment date, copied exactly: \texttt{2026-08-21}. This is supplied dated curation, not independently verified. It labels the problem partially solved and states that Madsen--Weiss settles the surface case while the extracted higher-dimensional generalization is not a general theorem. The background expressly requests independent verification of time-sensitive status and cited claims. The registry also supplies the statement-status label \texttt{corrected\_verified} and a claim that the formulation was checked against the official report; neither label nor claim constitutes independent verification in this draft.
\end{itemize}
\end{document}
